\documentclass[UTF8,zihao=5,fontset=fandol]{ctexart}

\usepackage[letterpaper,top=1in,bottom=1in,left=1.25in,right=1.25in]{geometry}
\usepackage{fontspec}
\usepackage{graphicx}
\usepackage{xcolor}
\usepackage{listings}
\usepackage{enumitem}
\usepackage{array}
\usepackage[most]{tcolorbox}
\usepackage[hidelinks]{hyperref}

\IfFontExistsTF{Times New Roman}
  {\setmainfont{Times New Roman}}
  {\setmainfont{TeX Gyre Termes}}

% ============================================================
% 模板固定文字（课程信息、教师信息、各段标题等）
% ============================================================
% 模板固定信息
% 本文件用于维护课程信息、教师信息、段落标题和说明文字。
\newcommand{\AttachmentText}{附件（四）}
\newcommand{\ReportTitle}{深\hspace{0.45em}圳\hspace{0.45em}大\hspace{0.45em}学\hspace{0.45em}实\hspace{0.45em}验\hspace{0.45em}报\hspace{0.45em}告}
\newcommand{\OfficeText}{教务处制}

\newcommand{\CourseLabel}{课程名称：}
\newcommand{\ExperimentLabel}{实验项目名称：}
\newcommand{\CollegeLabel}{学院：}
\newcommand{\MajorLabel}{专业：}
\newcommand{\InstructorLabel}{指导教师：}
\newcommand{\StudentLabel}{报告人：}
\newcommand{\ClassLabel}{班级：}
\newcommand{\ExperimentDateLabel}{实验时间：}
\newcommand{\SubmissionDateLabel}{实验报告提交时间：}

\newcommand{\CourseName}{微处理器与机器人}
\newcommand{\ExperimentName}{跑马灯点亮实验}
\newcommand{\CollegeName}{计算机与软件学院}
\newcommand{\InstructorName}{王东升}
% 报告人和学号组成一个整体，并沿封面的全局中心线居中。
\newcommand{\StudentLine}{%
  \makebox[2.0cm][c]{\StudentName}%
  \hspace{0.4cm}学号：\makebox[2.4cm][c]{\StudentID}%
}

% 各段标题与模板说明
\newcommand{\PurposeTitle}{一、实验目的与要求：}
\newcommand{\MethodTitle}{二、方法、步骤：}
\newcommand{\MethodHint}{（说明程序相关的算法原理或知识内容，程序设计的思路和方法，可以用流程图表述，程序主要数据结构的设计、主要函数之间的调用关系等）}
\newcommand{\ProcessTitle}{三、实验过程及内容：}
\newcommand{\ProcessHint}{（对程序代码进行说明和分析，越详细越好，代码排版要整齐，可读性要高）}
\newcommand{\ConclusionTitle}{四、实验结论：}
\newcommand{\ConclusionHint}{（提供运行结果，对结果进行探讨、分析、评价，并提出结论性意见和改进想法）}
\newcommand{\ReflectionTitle}{五、实验体会：}
\newcommand{\ReflectionHint}{（根据自己情况填写）}
\newcommand{\ReviewNote}{注：“指导教师批阅意见”栏请单独放置一页}
\newcommand{\ReviewTitle}{指导教师批阅意见：}
\newcommand{\GradeTitle}{成绩评定：}
\newcommand{\SignatureTitle}{指导教师签字：}
\newcommand{\RemarksTitle}{备注：}


% ============================================================
% 学生信息填写区
% 只需修改下面各命令第二对花括号中的内容。
% ============================================================
\newcommand{\MajorName}{}              % 专业
\newcommand{\StudentName}{}            % 姓名
\newcommand{\StudentID}{}              % 学号
\newcommand{\ClassName}{}              % 班级
\newcommand{\ExperimentDate}{年\quad 月\quad 日} % 实验时间
\newcommand{\SubmissionDate}{年\quad 月\quad 日}         % 实验报告提交时间

% ============================================================
% 模板排版设置（通常不需要修改）
% ============================================================
\pagestyle{empty}
\setlength{\parindent}{2em}
\setlength{\parskip}{0pt}
\setlength{\fboxsep}{5pt}
\setlength{\fboxrule}{0.4pt}
\setlist{nosep,leftmargin=2em}

\definecolor{codegray}{RGB}{245,245,245}
\lstset{
  basicstyle=\ttfamily\small,
  backgroundcolor=\color{codegray},
  frame=single,
  rulecolor=\color{black!30},
  numbers=left,
  numberstyle=\tiny\color{black!55},
  breaklines=true,
  columns=fullflexible,
  keepspaces=true,
  showstringspaces=false,
  tabsize=4
}

% 固定宽度的带下划线字段。
\newcommand{\reportfield}[2]{%
  \underline{\makebox[#1][c]{\strut #2}}%
}

% 封面各行左端对齐，下划线右端对齐，内容按整行的中心线对齐。
\newlength{\coverrowwidth}
\newlength{\coverlabelwidth}
\newlength{\coverfieldwidth}
\setlength{\coverrowwidth}{13.8cm}
\newcommand{\coverrow}[2]{%
  \settowidth{\coverlabelwidth}{#1}%
  \setlength{\coverfieldwidth}{\dimexpr\coverrowwidth-\coverlabelwidth-0.5em\relax}%
  \makebox[0pt][l]{#1\hspace{0.5em}\reportfield{\coverfieldwidth}{}}%
  \makebox[\coverrowwidth][c]{\strut #2}%
}

% 正文页的外框尺寸与原 Word 模板一致。
\newlength{\reportinnerheight}
\setlength{\reportinnerheight}{\dimexpr\textheight-2\fboxsep-2\fboxrule\relax}
\newsavebox{\reportpagebox}
\newenvironment{reportpage}
  {\begin{lrbox}{\reportpagebox}\begin{minipage}[t][\reportinnerheight][t]{\dimexpr\textwidth-2\fboxsep-2\fboxrule\relax}}
  {\end{minipage}\end{lrbox}\noindent\fbox{\usebox{\reportpagebox}}}

\newcommand{\reportheading}[2]{%
  \noindent\textbf{#1}%
  \if\relax\detokenize{#2}\relax\else #2\fi\par
}

% 正文内容框会根据内容自动增高，并在空间不足时自动分页。
\newtcolorbox{reportsection}[2]{
  enhanced jigsaw,
  breakable,
  colback=white,
  colframe=black,
  boxrule=\fboxrule,
  arc=0pt,
  outer arc=0pt,
  left=5pt,
  right=5pt,
  top=5pt,
  bottom=5pt,
  before skip=0pt,
  after skip=0pt,
  before upper={\reportheading{#1}{#2}}
}

\begin{document}

% ============================================================
% 第 1 页：封面
% ============================================================
\noindent \AttachmentText

\vspace{2.2em}
{\centering\bfseries\fontsize{22pt}{28pt}\selectfont \ReportTitle\par}

\vspace{2.8em}
{\bfseries\fontsize{14pt}{21pt}\selectfont
\setlength{\parindent}{0pt}
\centering
\coverrow{\CourseLabel}{\CourseName}\par\vspace{1.5em}
\coverrow{\ExperimentLabel}{\ExperimentName}\par\vspace{1.5em}
\coverrow{\CollegeLabel}{\CollegeName}\par\vspace{1.5em}
\coverrow{\MajorLabel}{\MajorName}\par\vspace{1.5em}
\coverrow{\InstructorLabel}{\InstructorName}\par\vspace{1.5em}
\coverrow{\StudentLabel}{\StudentLine}\par\vspace{1.5em}
\coverrow{\ClassLabel}{\ClassName}\par\vspace{1.5em}
\coverrow{\ExperimentDateLabel}{\ExperimentDate}\par\vspace{1.5em}
\coverrow{\SubmissionDateLabel}{\SubmissionDate}\par
}

\vfill
{\centering\bfseries\fontsize{12pt}{16pt}\selectfont \OfficeText\par}

\newpage

% ============================================================
% 报告正文：各部分按内容自动增高、自动分页
% ============================================================
\begin{reportsection}{\PurposeTitle}{}
  % ---------- “实验目的与要求”填写区开始 ----------
  跑马灯是嵌入式系统入门实验中常见的人机交互实例。本实验以 STM32F4 开发板为平台，将 LED、独立按键和蜂鸣器等外设组合起来，使学生理解微控制器如何采集外部输入、执行控制逻辑并驱动输出设备，为后续学习中断、定时器及复杂外设控制打下基础。

  \textbf{相关预备知识：}
  \begin{itemize}
    \item \textbf{GPIO 基础：}了解 STM32F4 GPIO 的输入、推挽输出、上拉和下拉等工作模式，能够根据开发板原理图确认 LED、按键和蜂鸣器所连接的端口、引脚及有效电平。
    \item \textbf{LED 驱动：}理解 LED 的限流和高、低电平点亮方式，掌握使用 \texttt{HAL\_GPIO\_WritePin()}、\texttt{HAL\_GPIO\_TogglePin()} 等函数改变 LED 状态。
    \item \textbf{按键检测：}掌握使用 \texttt{HAL\_GPIO\_ReadPin()} 读取按键电平的方法，理解机械按键抖动的产生原因，以及延时消抖、状态变化检测和单次触发的基本思路。
    \item \textbf{时序与状态控制：}理解阻塞式延时对程序响应速度的影响，能够使用循环、计时变量或状态机描述跑马灯的执行顺序和灯效切换过程。
    \item \textbf{蜂鸣器控制：}根据开发板所用蜂鸣器的类型和有效电平选择正确的驱动方式；按键触发后应控制鸣响持续时间，防止蜂鸣器因按键持续按下而连续发声。
    \item \textbf{PWM 基础：}了解 STM32 定时器、PWM 频率和占空比的含义，理解通过周期性改变占空比调节 LED 平均亮度并形成呼吸效果的原理。
  \end{itemize}

  \textbf{实验目的：}
  \begin{enumerate}[label=\arabic*)]
    \item 熟悉 STM32F4 开发板的基本结构，以及 STM32CubeIDE 中工程创建、引脚配置、代码生成、编译、下载和调试的基本流程。
    \item 理解 GPIO 输入与输出的工作原理，掌握使用 HAL 库读取按键状态、控制 LED 和驱动蜂鸣器的方法。
    \item 学习利用循环、延时、位操作及状态变量组织灯光变化，并通过函数封装提高程序的模块化程度和可读性。
    \item 掌握按键消抖和按键事件检测方法，理解“按下事件”与“持续按下状态”的区别，避免一次按键产生多次误触发。
    \item 了解定时器 PWM 的频率、占空比与 LED 亮度之间的关系，为实现呼吸灯及后续电机调速等功能建立基础。
  \end{enumerate}

  \textbf{实验要求：}使用 8 个 LED 实现 4 种具有明显差异的跑马灯特效；使用 4 个按键分别选择对应特效；每次有效按键按下时，蜂鸣器只发出一次提示音，并保证灯效能够稳定切换。程序应合理划分功能模块，使用含义明确的变量和注释，避免各外设控制之间相互干扰。选做部分使用 PWM 占空比渐变实现亮度变化平滑的呼吸灯效果。
  % ---------- “实验目的与要求”填写区结束 ----------
\end{reportsection}

\begin{reportsection}{\MethodTitle}{\MethodHint}
  % ---------- “方法、步骤”填写区开始 ----------
  % 请在这里填写内容。
  % ---------- “方法、步骤”填写区结束 ----------
\end{reportsection}

\begin{reportsection}{\ProcessTitle}{\ProcessHint}
  % ---------- “实验过程及内容”填写区开始 ----------
  % 请在这里填写内容。
  % ---------- “实验过程及内容”填写区结束 ----------
\end{reportsection}

\begin{reportsection}{\ConclusionTitle}{\ConclusionHint}
  % ---------- “实验结论”填写区开始 ----------
  % 请在这里填写内容。
  % ---------- “实验结论”填写区结束 ----------
\end{reportsection}

\begin{reportsection}{\ReflectionTitle}{\ReflectionHint}
  % ---------- “实验体会”填写区开始 ----------
  % 请在这里填写内容。
  % ---------- “实验体会”填写区结束 ----------

  \par\medskip\noindent \ReviewNote
\end{reportsection}

\clearpage

% ============================================================
% 第 5 页：教师批阅页（学生无需填写）
% ============================================================
\begin{reportpage}
  \begin{minipage}[t][0.56\reportinnerheight][t]{\linewidth}
    \noindent \ReviewTitle
  \end{minipage}
  \par\nointerlineskip\rule{\linewidth}{\fboxrule}\par
  \begin{minipage}[t][0.15\reportinnerheight][t]{\linewidth}
    \vspace{5pt}\noindent \GradeTitle
  \end{minipage}
  \par\nointerlineskip\rule{\linewidth}{\fboxrule}\par
  \begin{minipage}[t][0.16\reportinnerheight][t]{\linewidth}
    \vspace{5pt}
    \hfill \SignatureTitle\hspace{2em}\par
    \vspace{1.5em}
    \hfill 年\hspace{2em}月\hspace{2em}日\hspace{2em}
  \end{minipage}
  \par\nointerlineskip\rule{\linewidth}{\fboxrule}\par
  \vspace{5pt}\noindent \RemarksTitle
\end{reportpage}

\end{document}
